\section{Conclusiones y trabajo futuro}

La presente revisión sistemática permite afirmar categóricamente que la transición hacia Arquitecturas Orientadas a Eventos (EDA) y microservicios en el ecosistema IoT ha dejado de ser una exploración académica para convertirse en una **necesidad técnica impuesta por la física** de los datos modernos. La evidencia empírica analizada demuestra que las arquitecturas tradicionales, centralizadas y monolíticas, como la descrita en el caso de estudio de Tecgraf \cite{laigner2020monolithic}, han alcanzado un límite funcional insuperable. Ya no poseen la capacidad elástica ni la velocidad de conmutación necesarias para gestionar la concurrencia que requieren escenarios críticos actuales, como la prevención autónoma de colisiones espaciales mediante visión neuromórfica \cite{coretti2025neuromorphic} o el control de tráfico masivo en tiempo real \cite{perera2024traffic}.

\subsection{Síntesis de Hallazgos Técnicos}
Del análisis cruzado y exhaustivo de los 20 casos de estudio, derivamos cinco pilares fundamentales que definen el estado del arte actual y separan las implementaciones exitosas de los fracasos arquitectónicos:

\begin{enumerate}
    \item \textbf{La supremacía del Borde (Edge):} El modelo de "nube céntrica" se ha vuelto económicamente y técnicamente inviable para aplicaciones de latencia crítica. La implementación de filtrado inteligente en el borde, documentada extensamente por Baz et al. \cite{baz2019video} y Salamon et al. \cite{salamon2018gateway}, valida que procesar el dato "donde nace" no solo reduce la latencia a milisegundos, sino que disminuye drásticamente los costos operativos (OPEX) de ancho de banda al evitar la transmisión de terabytes de "ruido digital" irrelevante hacia los servidores centrales.
    
    \item \textbf{Eficiencia sobre Fuerza Bruta:} No siempre la solución es escalar verticalmente el hardware. Estudios como el de Kul y Tashiev \cite{kul2021microservices} demostraron empíricamente que modelos de IA optimizados y ligeros (como Tiny-YOLO) pueden superar en eficiencia operativa a modelos complejos si se integran correctamente en una arquitectura de flujo paralelo (Kafka Streams). La optimización de la topología del software rinde más frutos que la mera adición de potencia de cómputo.
    
    \item \textbf{El valor del Contexto Semántico:} La velocidad es inútil sin interpretación. La mejora del 54\% al 84\% en la detección de matrículas reportada por Perera et al. \cite{perera2024traffic} prueba que la lógica de negocio específica (algoritmos \textit{ad-hoc}) es tan crucial como el motor de IA subyacente. Sin una capa que entienda el contexto local y las reglas del dominio, los algoritmos genéricos fallan, generando datos basura a alta velocidad.
    
    \item \textbf{Estructuración Lógica de Datos:} La velocidad de red no compensa una mala organización. Shujaat et al. \cite{shujaat2021improvements} demostraron que en entornos IoT masivos, la latencia de recuperación de datos es el verdadero cuello de botella. Su propuesta de organizar los sensores mediante "Árboles de Acceso a Servicios" (SAT) en lugar de listas planas confirma que la estructura de datos es tan vital como el protocolo de transmisión.
    
    \item \textbf{De la Detección a la Respuesta:} Finalmente, identificamos que detectar un evento es solo la mitad de la ecuación. Como argumentan Ward et al. \cite{ward2005response}, un sistema verdaderamente resiliente debe implementar una "Arquitectura de Respuesta Dirigida por Eventos" (EDRA). No basta con saber que algo ocurrió; el sistema debe tener la inteligencia contextual para decidir automáticamente qué acción tomar, cerrando el ciclo entre la percepción digital y la acción física.
\end{enumerate}

\subsection{Implicaciones Estratégicas para la Industria}
Basándonos en los patrones de éxito y fracaso identificados en la literatura, ofrecemos las siguientes directrices estratégicas para arquitectos de software y líderes técnicos que planean una migración:

\begin{itemize}
    \item \textbf{Gestión de la Complejidad Híbrida:} La asincronía introduce una complejidad de depuración significativa. Para procesos que requieren consistencia transaccional estricta (ACID), como pagos o multas legales, se recomienda mantener un control orquestado centralizado, tal como sugiere Adekunle \cite{adekunle2024financial}. La industria debe moverse hacia modelos híbridos donde la coreografía de eventos maneja la ingesta masiva, pero la orquestación BPM maneja el estado del negocio.
    
    \item \textbf{Invertir en la Capa Semántica:} Antes de conectar sensores masivamente, es imperativo definir una ontología común. Como advierten Phuttharak y Loke \cite{phuttharak2023smartcity} y Xu et al. \cite{xu2016provisioning}, la falta de estandarización crea "silos de datos" que, aunque rápidos, son incapaces de cooperar entre sí, anulando el propósito de una *Smart City*. La integración debe ser semántica, no solo sintáctica.
    
    \item \textbf{Evolución Incremental (EDSOA):} Para sistemas legados que no pueden ser reescritos desde cero (Brownfield), la adopción progresiva de una Arquitectura Orientada a Servicios Dirigida por Eventos (EDSOA), propuesta por Shi et al. \cite{shi2015execution}, ofrece un camino intermedio viable. Esto permite envolver servicios antiguos con interfaces de eventos, inyectando reactividad sin el riesgo de una reingeniería total ("Big Bang").
\end{itemize}

\subsection{Impacto Socioeconómico}
Más allá del código, esta revisión destaca el impacto tangible de estas arquitecturas en la sociedad. La capacidad de inferir el estado de una ambulancia en tiempo real mediante CEP, como demostraron Bruns y Dunkel \cite{bruns2014fusion}, no es solo una mejora de eficiencia; se traduce directamente en vidas salvadas al reducir los tiempos de despacho. De igual manera, las plataformas de decisión para ciudades inteligentes propuestas por Visnjic et al. \cite{visnjic2021decision} sugieren que la adopción de EDA es un habilitador clave para la sostenibilidad urbana, permitiendo una gestión de recursos basada en la demanda real y no en estimaciones estáticas.

\subsection{Desafíos Abiertos y Trabajo Futuro}
A pesar de los avances, la modernización no está exenta de riesgos críticos. Identificamos que la **observabilidad** sigue siendo el talón de Aquiles de estas arquitecturas; sin herramientas de trazabilidad distribuida, los sistemas EDA se convierten en "cajas negras" imposibles de auditar en caso de fallo, un riesgo inaceptable expuesto por Lazzari y Farias \cite{lazzari2023hidden}.

El trabajo a futuro en esta línea de investigación debe concentrarse obsesivamente en tres frentes:
\begin{enumerate}
    \item \textbf{Monitoreo Semántico:} Desarrollar herramientas que permitan rastrear no solo si un mensaje llegó, sino si su contenido (payload) mantuvo su integridad lógica a través de múltiples transformaciones asíncronas en microservicios dispares.
    \item \textbf{Eficiencia Energética en Protocolos:} Seguir la línea de Park et al. \cite{park2021cloud} para optimizar el consumo de energía en la transmisión de eventos. A medida que el IoT escala, la huella de carbono de la transmisión de datos se vuelve un factor limitante que requiere protocolos más ligeros que HTTP o TCP estándar.
    \item \textbf{Interoperabilidad Plug-and-Play:} Establecer estándares de interoperabilidad ligera que no requieran ontologías pesadas, facilitando que dispositivos de diferentes fabricantes colaboren en el borde mediante protocolos ágiles como MQTT, eliminando la fricción de integración actual.
\end{enumerate}

En definitiva, la arquitectura distribuida ha dejado de ser una opción para convertirse en el estándar de supervivencia. Mientras no resolvamos la dificultad de observar y depurar estos sistemas, su adopción masiva fuera de los gigantes tecnológicos seguirá siendo una apuesta arriesgada, donde se intercambia estabilidad por velocidad.